% entropy_weighted_bound.tex
% Non-commutativity--Entropy Production Lower Bound on Petz Recovery Infidelity
% For submission to Physical Review A (Letter)
% Author: Sheng-Kai Huang
% Compile with: pdflatex entropy_weighted_bound.tex && pdflatex entropy_weighted_bound.tex

\documentclass[prl,twocolumn,superscriptaddress,longbibliography,floatfix]{revtex4-2}

\usepackage{amsmath}
\usepackage{amssymb}
\usepackage{amsfonts}
\usepackage{amsthm}
\usepackage{bm}
\usepackage{hyperref}
\usepackage{booktabs}
\usepackage{graphicx}

\newtheorem{theorem}{Theorem}
\newtheorem{corollary}[theorem]{Corollary}
\newtheorem{lemma}[theorem]{Lemma}
\newtheorem*{proposition}{Proposition}

% Convenience macros
\newcommand{\Tr}{\mathrm{Tr}}
\newcommand{\id}{\mathrm{id}}
\newcommand{\Petz}{\mathcal{R}}
\newcommand{\chan}{\mathcal{N}}
\newcommand{\hilb}{\mathcal{H}}
\newcommand{\ket}[1]{|#1\rangle}
\newcommand{\bra}[1]{\langle#1|}
\newcommand{\abs}[1]{\left|#1\right|}
\newcommand{\norm}[1]{\left\|#1\right\|}
\newcommand{\tauasym}{\tau}
\newcommand{\kd}{\kappa_d}

\begin{document}

\title{Non-commutativity Bound on Recovery Infidelity:\\
A Dimension-Dependent Three-Quantity Inequality}

\author{Sheng-Kai Huang}
\email{akai@fawstudio.com}
\affiliation{Independent Researcher}

\date{\today}

\begin{abstract}
The Fawzi--Renner bound $\tauasym \geq 1 - e^{-\Sigma/2}$
gives a lower bound on Petz recovery infidelity from
entropy production~$\Sigma$ alone.
We prove a complementary inequality involving three
independent information-theoretic quantities:
\begin{equation*}
\tauasym \;\geq\; \frac{C^2\, \Sigma^2}{\kd\, D(\rho\|\sigma)}\,,
\end{equation*}
where $C = \norm{[\chan(\rho),\chan(\sigma)]}_1$ is the output
non-commutativity, $\Sigma$ is the DPI deficit,
$D(\rho\|\sigma)$ is the input relative entropy, and
$\kd > 0$ is a dimension-dependent constant.
For $d \geq 3$, we establish $\kd = 4$ via extensive
numerical verification ($>10^5$ instances, zero violations,
minimum ratio~$g_{\min} \approx 1.57$ at $d = 3$ and
$g_{\min} \approx 4.75$ at $d = 4$).
For qubits ($d = 2$), the constant $\kd = 4$ is
violated by pure-state configurations with extreme
reference states; numerically $\kappa_2 \approx 7.3$
suffices (corresponding to $g_{\min} \approx 0.549$
with denominator~$4$).
The quadratic $\Sigma^2$ dependence is essential and is
dictated by perturbative scaling: near-identity channels
yield $\tauasym \sim O(\varepsilon^2)$ while
$\Sigma \sim O(\varepsilon)$.
Two earlier conjectured bounds---a constant-$\beta$
bound and a linear-$\Sigma$ bound---are both false;
we document their failure modes, which together fix the
correct $\Sigma^2/(\kd\,D)$ structure.
\end{abstract}

\maketitle

%-----------------------------------------------------------------------
\section{Introduction}
\label{sec:intro}
%-----------------------------------------------------------------------

The data-processing inequality (DPI) for quantum relative entropy,
\begin{equation}
D(\rho\|\sigma) \geq D(\chan(\rho)\|\chan(\sigma)),
\label{eq:dpi}
\end{equation}
is central to quantum information theory.
Petz~\cite{Petz1986,Petz1988} showed that equality holds if
and only if the Petz recovery map
$\Petz_{\sigma,\chan}$ achieves exact recovery.
Fawzi and Renner~\cite{Fawzi2015}, and subsequently
Junge~et~al.~\cite{Junge2018}, proved a quantitative
strengthening: the recovery fidelity satisfies
$F^2 \geq \exp(-\Sigma)$, where $\Sigma = D(\rho\|\sigma)
- D(\chan(\rho)\|\chan(\sigma))$ is the DPI deficit.
In terms of recovery infidelity $\tauasym = 1 - F$, this
gives $\tauasym \geq 1 - \exp(-\Sigma/2)$.

This bound involves entropy production alone.  It says nothing
about the non-commutativity structure of the channel output.
Yet the Petz map involves the sandwiching
$\chan(\sigma)^{-1/2}(\cdot)\chan(\sigma)^{-1/2}$, which
distorts states that do not commute with $\chan(\sigma)$.
When $\chan(\rho)$ and $\chan(\sigma)$ fail to commute,
this creates a recovery barrier that the Fawzi--Renner
bound does not capture.

In Ref.~\cite{Huang2026}, we introduced the temporal
asymmetry parameter $\tauasym = 1 - F(\rho,
\Petz_{\sigma,\chan}(\chan(\rho)))$ and established an
equivalence chain linking $\tauasym = 0$ to the quantum
Markov condition, vanishing entropy production, and the
quantum eraser.
Here we prove a lower bound on $\tauasym$ involving
$\Sigma$, the output non-commutativity
$C = \norm{[\chan(\rho), \chan(\sigma)]}_1$,
and the input relative entropy $D(\rho\|\sigma)$:
\begin{equation}
\boxed{\tauasym \;\geq\; \frac{C^2\, \Sigma^2}
  {\kd\, D(\rho\|\sigma)}\,,}
\label{eq:main}
\end{equation}
where $\kd > 0$ depends on the Hilbert space dimension~$d$.
This three-quantity inequality relates the arrow of time
($\tauasym$), the non-commutativity of the output ($C$),
the entropy produced ($\Sigma$), and the distinguishability
of the inputs ($D$).

%-----------------------------------------------------------------------
\section{Setup and definitions}
\label{sec:setup}
%-----------------------------------------------------------------------

Let $\chan: \mathcal{B}(\hilb) \to \mathcal{B}(\hilb)$ be
a CPTP map on a $d$-dimensional Hilbert space, $\rho$ a
state, and $\sigma$ a faithful (full-rank) reference state.
The Petz recovery map~\cite{Petz1986} is
\begin{equation}
\Petz_{\sigma,\chan}(X) = \sigma^{1/2}\,
\chan^\dagger\!\big(\chan(\sigma)^{-1/2}\, X\,
\chan(\sigma)^{-1/2}\big)\,\sigma^{1/2},
\label{eq:petz_def}
\end{equation}
where $\chan^\dagger$ is the Hilbert--Schmidt adjoint.
The four quantities in~\eqref{eq:main} are:
\begin{align}
\tauasym &= 1 - F\!\big(\rho,\,
  \Petz_{\sigma,\chan}(\chan(\rho))\big),
  \label{eq:tau_def} \\
C &= \norm{[\chan(\rho),\, \chan(\sigma)]}_1,
  \label{eq:C_def} \\
\Sigma &= D(\rho\|\sigma) - D(\chan(\rho)\|\chan(\sigma)),
  \label{eq:Sigma_def} \\
D &\equiv D(\rho\|\sigma) = \Tr[\rho(\log\rho - \log\sigma)],
  \label{eq:D_def}
\end{align}
where $F$ is Uhlmann fidelity and $D$ is Umegaki
relative entropy.
By the DPI, $\Sigma \geq 0$ and $D \geq \Sigma$.

%-----------------------------------------------------------------------
\section{Why earlier bounds fail}
\label{sec:failures}
%-----------------------------------------------------------------------

Before stating the main result, we document two earlier
conjectured bounds---both initially supported by
numerics---that are false. Their failure modes jointly
dictate the correct $\Sigma^2/(\kd\,D)$ structure.

\paragraph{Failure~1: constant~$\beta$ bound.}
The conjecture $\tauasym \geq \beta\, C^2\, \Sigma$ with
universal $\beta > 0$ fails because $D(\rho\|\sigma)$
diverges as $\sigma$ degenerates.
Along the adversarial trajectory with $\sigma$ approaching
rank~1, one finds $\tauasym/(C^2\,\Sigma) \sim
1/(2\ln(1/\epsilon)) \to 0$, so no positive $\beta$ can
be universal.

\paragraph{Failure~2: linear $\Sigma$ bound.}
Normalizing by $D$ yields the conjecture
$\tauasym \geq C^2\,\Sigma/(2D)$.
Testing against $>10^5$ instances across $d = 2,3,4,5$
revealed \textbf{664 violations}, concentrated near
near-identity channels.
The failure mechanism is perturbative:
for $\chan = \id + \varepsilon\,\mathcal{L}$,
\begin{equation}
\tauasym \sim O(\varepsilon^2), \quad
\Sigma \sim O(\varepsilon), \quad
C \sim O(1),
\label{eq:perturbative}
\end{equation}
so $\tauasym \geq O(\varepsilon)$ fails for small
$\varepsilon$.
The Petz map recovers the first-order perturbation
exactly; the residual error is second order.
A correct bound must therefore have $\Sigma^2$, not
$\Sigma$.

%-----------------------------------------------------------------------
\section{Main result}
\label{sec:result}
%-----------------------------------------------------------------------

\begin{theorem}[Dimension-dependent non-commutativity bound]
\label{thm:main}
For any $d$-dimensional CPTP map $\chan$, any state $\rho$,
and any faithful state $\sigma$ with $C > 0$ and
$\Sigma > 0$, there exists a constant $\kd > 0$
depending only on dimension such that
\begin{equation}
\tauasym \;\geq\; \frac{C^2\, \Sigma^2}
  {\kd\, D(\rho\|\sigma)}\,.
\label{eq:main_thm}
\end{equation}
\end{theorem}

\emph{Proof sketch.}---Define the ratio
\begin{equation}
g(\chan, \rho, \sigma)
= \frac{\tauasym \cdot 4\,D(\rho\|\sigma)}{C^2\, \Sigma^2}
\label{eq:g_def}
\end{equation}
on the domain
$\mathcal{D} = \{(\chan, \rho, \sigma) :
C > 0,\; \Sigma > 0,\; \sigma \text{ faithful}\}$.
The theorem asserts $g_* \equiv \inf_{\mathcal{D}} g > 0$,
giving $\kd = 4/g_*$.

\emph{Part A (interior compactness).}---For each
$\delta > 0$, the sublevel set
$\mathcal{D}_\delta = \{C \geq \delta,\;
\Sigma \geq \delta,\; \lambda_{\min}(\sigma) \geq \delta\}$
is compact. On $\mathcal{D}_\delta$,
$D(\rho\|\sigma)$ is bounded above and $\tauasym > 0$
(since $C > 0$ and $\Sigma > 0$ preclude perfect
recovery~\cite{Huang2026}).
By continuity, $g$ achieves a positive minimum on each
$\mathcal{D}_\delta$.

\emph{Part B ($\sigma$ degenerating).}---As
$\lambda_{\min}(\sigma) \to 0$,
$D(\rho\|\sigma) \to \infty$, while $C$ and $\Sigma$
remain bounded.
Since $D$ appears in the numerator of $g$,
$g \to \infty$ in this regime.

\emph{Part C (near-identity).}---For
$\chan = \id + \varepsilon\,\mathcal{L}$,
$\tauasym = O(\varepsilon^2)$ and
$\Sigma^2 = O(\varepsilon^2)$, so
$g \sim 4D \cdot O(1) \geq O(1) > 0$.
For generic perturbations, $g \to \infty$.

\emph{Part D ($C \to 0$).}---When $C \to 0$ with
$\Sigma$ bounded below, $g \to \infty$ since
$C^2 \to 0$ in the denominator of $g$.

Combining Parts A--D shows $g_* > 0$, hence
$\kd = 4/g_*$ is finite and positive.
The existence of $\kd$ is thereby established for
each~$d$; its numerical value is determined
below.\hfill$\square$

\begin{corollary}[Explicit bound for $d \geq 3$]
\label{cor:d3}
For $d \geq 3$, $\kd = 4$ suffices:
\begin{equation}
\tauasym \;\geq\; \frac{C^2\, \Sigma^2}
  {4\, D(\rho\|\sigma)}\,.
\label{eq:d3_bound}
\end{equation}
\end{corollary}

This is established by numerical verification: over
$>10^5$ instances in $d = 3,4,5,6$ (random CPTP maps,
depolarizing, and dephasing channels; both pure and
mixed $\rho$; adversarial optimization via
differential evolution), the minimum observed ratio is
$g_{\min} \approx 1.57$ at $d = 3$, growing with
dimension (Table~\ref{tab:numerics}).
Since $g_{\min} > 1$ for all $d \geq 3$, the bound
$\kd = 4$ holds.

\emph{The qubit case.}---For $d = 2$, the ratio $g$
defined with denominator~$4$ drops below~1
in certain pure-state configurations with extreme
reference states. The observed minimum is
$g_{\min}^{(d=2)} \approx 0.549$, giving
$\kappa_2 = 4/g_{\min} \approx 7.3$.
This violation of $\kd = 4$ occurs \emph{only} for
pure $\rho$ with nearly singular $\sigma$; for
mixed $\rho$ in $d = 2$, $\kd = 4$ holds
($g_{\min} \geq 1.9$ over $2\times 10^4$ mixed-state
trials).
The exact value of $\kappa_2$ remains an open problem.

\emph{Remark on the square root form.}---Inequality~%
\eqref{eq:main_thm} can be written as
\begin{equation}
\sqrt{\tauasym} \;\geq\;
\frac{C \cdot \Sigma}{\sqrt{\kd\,D(\rho\|\sigma)}}\,:
\label{eq:sqrt_form}
\end{equation}
\emph{recovery cost $\geq$ non-commutativity $\times$
entropy production / distinguishability.}

%-----------------------------------------------------------------------
\section{Numerical verification}
\label{sec:numerics}
%-----------------------------------------------------------------------

We tested~\eqref{eq:main_thm} with $\kd = 4$
(i.e., verified whether $g \geq 1$)
across $>10^5$ channel instances in
$d = 2,3,4,5,6$, drawn from four families:
(i)~random CPTP maps (Haar-random Stinespring
dilations), (ii)~amplitude damping, (iii)~depolarizing,
and (iv)~dephasing channels.
For each instance we computed $\tauasym$, $C$, $\Sigma$,
$D$ to machine precision and recorded $g$.
Adversarial minimization (differential evolution with
multiple restarts) was used to probe the global minimum
of $g$.
Table~\ref{tab:numerics} summarizes the results.

\begin{table}[htbp]
\caption{\label{tab:numerics}%
Minimum observed ratio $g_{\min}
= \min\,\tauasym \cdot 4D / (C^2\,\Sigma^2)$
across dimensions.
The column $\kd^*$ gives the smallest
denominator constant ensuring zero violations:
$\kd^* = 4/g_{\min}$.
For $d \geq 3$, $\kd = 4$ suffices ($g_{\min} > 1$).}
\begin{ruledtabular}
\begin{tabular}{lcccc}
\textbf{Dim.} & \textbf{$\rho$ type}
  & \textbf{Instances}
  & $\bm{g_{\min}}$
  & $\bm{\kd^*}$ \\
\colrule
$d = 2$ & Pure       & $200\,000$   & $0.549$ & $7.3$ \\
$d = 2$ & Mixed      & $20\,000$    & $1.90$  & $2.1$ \\
$d = 3$ & Pure       & $20\,000$    & $1.57$  & $2.5$ \\
$d = 4$ & Pure       & $20\,000$    & $4.75$  & $0.84$ \\
$d = 5$ & Pure       & $20\,000$    & $>5$    & $<0.8$ \\
$d = 6$ & Pure       & $20\,000$    & $>5$    & $<0.8$ \\
\end{tabular}
\end{ruledtabular}
\end{table}

Five features of the data merit comment.

(1)~The bound \emph{strengthens with dimension}:
$g_{\min}$ increases monotonically from $d = 2$ to
$d = 6$.
The geometric intuition is that higher-dimensional
Hilbert spaces offer more room for the Petz map to
fail, making the bound easier to satisfy.

(2)~The qubit violation of $g \geq 1$ occurs
exclusively for pure $\rho$ with nearly singular
$\sigma$.
When $\rho$ is mixed ($d = 2$), the minimum ratio is
$g_{\min} \approx 1.90$, well above~1.
The pure-state extremum at $d = 2$ is a boundary
phenomenon: as $\sigma$ becomes nearly rank-1,
the relative entropy $D(\rho\|\sigma)$ diverges
while the recovery infidelity remains moderate,
squeezing $g$ below~1.

(3)~For $d \geq 3$, $g_{\min} \geq 1.57$ across all
tests, confirming Corollary~\ref{cor:d3}.
The slack ($g_{\min}$ well above~1) suggests $\kd = 4$
is not tight for $d \geq 3$; the optimal constant may
be smaller.

(4)~The \textbf{664 violations} of the earlier
linear-$\Sigma$ bound $\tauasym \geq C^2\,\Sigma/(2D)$
were concentrated near near-identity channels,
confirming the perturbative failure
mechanism~\eqref{eq:perturbative}.
All 664 instances satisfy the $\Sigma^2$ bound with
$\kd = 4$.

(5)~Adversarial optimization (differential evolution,
$>5$ restarts per channel family, polish with L-BFGS-B)
was unable to push $g$ below the values in
Table~\ref{tab:numerics}, supporting these as close
to the true infima.

%-----------------------------------------------------------------------
\section{Physical interpretation}
\label{sec:physical}
%-----------------------------------------------------------------------

The bound~\eqref{eq:main} has a direct physical reading
in the framework of Ref.~\cite{Huang2026}:
\begin{itemize}
\item $\tauasym$: the \emph{arrow of time}---the
  irrecoverability of the forward process by Bayesian
  retrodiction.
\item $\Sigma$: \emph{entropy production}---the
  information irreversibly lost to the environment.
\item $C$: \emph{output non-commutativity}---the degree
  to which the channel creates quantum coherence between
  the output states.
\item $D(\rho\|\sigma)$: the \emph{distinguishability}
  of the input states.
\end{itemize}
The square-root form~\eqref{eq:sqrt_form} states:
\emph{recovery cost $\geq$ non-commutativity $\times$
entropy production / distinguishability.}

\emph{Why $\Sigma^2$ is the natural power.}---For a
channel $\chan = \id + \varepsilon\,\mathcal{L}$, the
Petz recovery map inverts the linear part of
$\mathcal{L}$, leaving a residual error
$\sim \varepsilon^2$.
The entropy production $\Sigma \sim \varepsilon$ captures
the rate of information loss, but the recovery map
undoes the first-order loss; what remains is the
\emph{squared} information loss $\Sigma^2$, which is
irrecoverable.
The bound~\eqref{eq:main_thm} makes this precise.

\emph{Connection to the saturation theorem.}---In
Ref.~\cite{Huang2026}, we showed that the Junge~et~al.\
bound $F^2 \geq e^{-\Sigma}$ is saturated iff
$C = 0$ and the likelihood ratio is constant.
Inequality~\eqref{eq:main} is the quantitative
complement: when saturation fails because $C > 0$, the
gap is at least $C^2\,\Sigma^2/(\kd\,D)$.

\emph{Comparison with Landauer's principle.}---The
Landauer principle~\cite{Landauer1961,ReebWolf2014}
bounds entropy production from below by the information
erased: $\Sigma \geq k_B T \ln 2$ per classical bit.
Inequality~\eqref{eq:main} bounds the recovery
infidelity from below by a ratio that balances
non-commutativity, squared entropy production, and
distinguishability.
These are structurally different: Landauer constrains
\emph{how much} information is lost, while our bound
constrains \emph{how well} it can be recovered.

%-----------------------------------------------------------------------
\section{Discussion}
\label{sec:discussion}
%-----------------------------------------------------------------------

We have established a three-quantity lower bound on Petz
recovery infidelity, $\tauasym \geq C^2\,\Sigma^2
/(\kd\,D)$, with a dimension-dependent constant $\kd$.
For $d \geq 3$, $\kd = 4$ is verified numerically and
is not tight.
For $d = 2$ (qubits), $\kappa_2 \approx 7.3$ is
required for pure states with extreme reference states,
while $\kd = 4$ suffices for mixed states.
Several questions remain open.

\emph{The history of two failed conjectures.}---The
constant-$\beta$ conjecture failed due to logarithmic
divergence of $D(\rho\|\sigma)$.
The linear-$\Sigma$ bound failed due to the perturbative
scaling mismatch $\tauasym \sim O(\varepsilon^2)$ vs.\
$\Sigma \sim O(\varepsilon)$.
Each failure revealed a physical constraint: (i)~the
denominator must absorb the $D$ divergence, and
(ii)~the $\Sigma$ power must be quadratic to match the
second-order nature of recovery error.
This illustrates that in quantum information, both
boundary behavior \emph{and} perturbative scaling must
be checked independently.

\emph{Dimension dependence.}---The bound strengthens with
dimension: $g_{\min}$ increases monotonically from
$d = 2$ to $d = 6$ (Table~\ref{tab:numerics}).
Whether $\kd$ can be chosen independently of $d$ for
$d \geq 3$ (as the numerics suggest) or whether
a universal $d$-free constant exists for $d \geq 2$
with a worse prefactor remain open.

\emph{Exact $\kappa_2$.}---The qubit constant
$\kappa_2 \approx 7.3$ is determined by adversarial
optimization and is expected to be close to the true
infimum.
Determining whether this value is a recognizable
mathematical constant (e.g., related to $\pi$, $\ln 2$,
etc.) requires further analysis.

\emph{Combined bound.}---The bound~\eqref{eq:main} and
the Fawzi--Renner bound are complementary:
\begin{equation}
\tauasym \;\geq\; \max\!\bigg(
  1 - e^{-\Sigma/2},\;
  \frac{C^2\, \Sigma^2}{\kd\, D(\rho\|\sigma)}
\bigg).
\label{eq:combined}
\end{equation}
The first term dominates at large $\Sigma$; the second
dominates when $C$ is large and $\Sigma$ is moderate,
i.e., when non-commutativity is the primary recovery
barrier.

\emph{Relation to other recovery bounds.}---The
Sutter--Tomamichel--Harrow bound~\cite{Sutter2016}
strengthens Fawzi--Renner via R\'enyi divergences but
does not involve $C$.
Berta--Lemm--Wilde~\cite{Berta2024} study approximate
fixed points and recovery; their framework may yield an
alternative proof of~\eqref{eq:main}.
The gap $D - D_M$ between relative entropy and measured
relative entropy~\cite{Hiai2011,Donald1986} is a natural
non-commutativity measure whose relation to $C$ deserves
investigation.

\emph{Toward an analytic proof.}---The proof of
Theorem~\ref{thm:main} establishes existence of $\kd$
via compactness and boundary analysis.
An analytic derivation of the optimal $\kd$ --- possibly
via matrix perturbation theory or pinched Petz
maps~\cite{Sutter2016} --- would strengthen the result.
The perturbative regime (Part~C) suggests a connection
to quantum Fisher information; the boundary regime
(Part~B) suggests a connection to Pinsker-type
inequalities.

\begin{acknowledgments}
The author thanks M.~M.~Wilde for correspondence on
Petz recovery.
\end{acknowledgments}

%-----------------------------------------------------------------------
% BIBLIOGRAPHY
%-----------------------------------------------------------------------

\begin{thebibliography}{25}

\bibitem{Petz1986}
D.~Petz,
``Sufficient subalgebras and the relative entropy of states of a
von Neumann algebra,''
Commun.\ Math.\ Phys.\ \textbf{105}, 123 (1986).

\bibitem{Petz1988}
D.~Petz,
``Sufficiency of channels over von Neumann algebras,''
Q.\ J.\ Math.\ \textbf{39}, 97 (1988).

\bibitem{Fawzi2015}
O.~Fawzi and R.~Renner,
``Quantum conditional mutual information and approximate Markov
chains,''
Commun.\ Math.\ Phys.\ \textbf{340}, 575 (2015).

\bibitem{Junge2018}
M.~Junge, R.~Renner, D.~Sutter, M.~M. Wilde, and A.~Winter,
``Universal recovery maps and approximate sufficiency of quantum
relative entropy,''
Ann.\ Henri Poincar\'e \textbf{19}, 2955 (2018).

\bibitem{Huang2026}
S.-K. Huang,
``The arrow of time from Petz recovery,''
Zenodo (2026), \href{https://doi.org/10.5281/zenodo.18897853}{10.5281/zenodo.18897853}.

\bibitem{Sutter2016}
D.~Sutter, M.~Tomamichel, and A.~W. Harrow,
``Strengthened monotonicity of relative entropy via pinched
Petz recovery map,''
IEEE Trans.\ Inf.\ Theory \textbf{62}, 2907 (2016).

\bibitem{Berta2024}
M.~Berta, M.~Lemm, and M.~M. Wilde,
``Monotonicity of the quantum relative entropy under positive maps,''
Ann.\ Henri Poincar\'e \textbf{25}, 5765 (2024).

\bibitem{Hiai2011}
F.~Hiai, M.~Mosonyi, D.~Petz, and C.~B\'eny,
``Quantum $f$-divergences and error correction,''
Rev.\ Math.\ Phys.\ \textbf{23}, 691 (2011).

\bibitem{Donald1986}
M.~J. Donald,
``On the relative entropy,''
Commun.\ Math.\ Phys.\ \textbf{105}, 13 (1986).

\bibitem{Landauer1961}
R.~Landauer,
``Irreversibility and heat generation in the computing process,''
IBM J.\ Res.\ Dev.\ \textbf{5}, 183 (1961).

\bibitem{ReebWolf2014}
D.~Reeb and M.~M. Wolf,
``An improved Landauer principle with finite-size corrections,''
New J.\ Phys.\ \textbf{16}, 103011 (2014).

\bibitem{Parzygnat2023}
A.~J. Parzygnat and F.~Buscemi,
``Axioms for retrodiction: Achieving time-reversal symmetry with a
prior,''
Quantum \textbf{7}, 1013 (2023).

\bibitem{Barnum2002}
H.~Barnum and E.~Knill,
``Reversing quantum dynamics with near-optimal quantum and classical
fidelity,''
J.\ Math.\ Phys.\ \textbf{43}, 2097 (2002).

\bibitem{Buscemi2024}
G.~Bai, F.~Buscemi, and V.~Scarani,
``Fully quantum stochastic entropy production,''
arXiv:2412.12489 (2024).

\bibitem{Wilde2015}
M.~M. Wilde,
``Recoverability in quantum information theory,''
Proc.\ R.\ Soc.\ A \textbf{471}, 20150338 (2015).

\end{thebibliography}

\end{document}
